%%%%%%%%%%%%%%%%
%%% deut 10 %%%
%%%%%%%%%%%%%%%%

\Chapter{10}

\Verse{1}
{
	\hDtr1{בָּעֵת הַהִוא אָמַר יְהֹוָה אֵלַי פְּסל לְךָ שְׁנֵי לוּחֹת אֲבָנִים כָּרִאשֹׁנִים וַעֲלֵה אֵלַי הָהָרָה וְעָשִׂיתָ לְּךָ אֲרוֹן עֵץ׃}
}
{
	\eDtr1{
		10 “At that time YHWH said to me, ‘Carve two tablets of
		stones like the first ones, and come up to me at the
		mountain. And you shall make an ark of wood.
	}
}
{
%	\footnote{
%		you shall make an ark of wood. It says “you shall make,” but in Exodus
%		it was Bezalel, not Moses, who made the ark. Rashi and Ramban and others
%		therefore say that there must have been two arks. In current critical
%		scholarship, the apparent contradiction would generally be taken to be a
%		result of the fact that this text in Deuteronomy and the text in Exodus
%		were written by two different authors. But, even without raising such
%		solutions that are arrived at through critical approaches, we can
%		understand Moses’ words in Deuteronomy to be a brief account of what
%		happened at Sinai, and so he speaks of himself as making the ark when he
%		means that he directed Bezalel to do it. Likewise, he speaks of the ark
%		as if it were finished and ready when he came down from the mountain
%		when he knows that actually some time passed before the ark was made. J.
%		H. Hertz proposed such an understanding as an alternative to Rashi’s and
%		Ramban’s view even though Hertz harshly rejected literary-critical
%		scholarship. This is important because an almost fundamentalist view of
%		Rashi is growing at present in some communities, wherein it is
%		practically heresy to question Rashi. This gives Rashi a status that I
%		do not think he would have wanted for himself. And certainly the
%		commentators of the generations that followed him, including his own
%		grandson Rashbam, were prepared to question his comments. We respect and
%		admire the great commentators of past generations, but we should never
%		forget to distinguish between Torah and commentary. The goal of refining
%		our knowledge of the Torah and shedding new light from it in each new
%		age is never-ending; and our commentaries—including my own—are always
%		subject to criticism and improvement.
%	}
}

\Verse{2}
{
	\hDtr1{וְאֶכְתֹּב עַל הַלֻּחֹת אֶת הַדְּבָרִים אֲשֶׁר הָיוּ עַל הַלֻּחֹת הָרִאשֹׁנִים אֲשֶׁר שִׁבַּרְתָּ וְשַׂמְתָּם בָּאָרוֹן׃}
}
{
	\eDtr1{
		And I’ll write on the tablets the words that were on the
		first tablets, which you shattered, and you shall set them
		in the ark.’
	}
}
{
}

\Verse{3}
{
	\hDtr1{וָאַעַשׂ אֲרוֹן עֲצֵי שִׁטִּים וָאֶפְסֹל שְׁנֵי לֻחֹת אֲבָנִים כָּרִאשֹׁנִים וָאַעַל הָהָרָה וּשְׁנֵי הַלֻּחֹת בְּיָדִי׃}
}
{
	\eDtr1{
		“And I made an ark of acacia wood, and I carved two
		tablets of stones like the first ones, and I went up the
		mountain with the two tablets in my hands.
	}
}
{
}

\Verse{4}
{
	\hDtr1{וַיִּכְתֹּב עַל הַלֻּחֹת כַּמִּכְתָּב הָרִאשׁוֹן אֵת עֲשֶׂרֶת הַדְּבָרִים אֲשֶׁר דִּבֶּר יְהֹוָה אֲלֵיכֶם בָּהָר מִתּוֹךְ הָאֵשׁ בְּיוֹם הַקָּהָל וַיִּתְּנֵם יְהֹוָה אֵלָי׃}
}
{
	\eDtr1{
		And He wrote on the tablets like the first writing: the
		Ten Commandments that YHWH spoke to you at the mountain
		from inside the fire in the day of the assembly, and YHWH
		gave them to me.
	}
}
{
%	\footnote{
%		the Ten Commandments. Hebrew ‘[image][image]eret hadd[image]b[image]rîm,
%		literally, the ten things. This is what the Ten Commandments are called
%		in Hebrew.
%	}
}

\Verse{5}
{
	\hDtr1{וָאֵפֶן וָאֵרֵד מִן הָהָר וָאָשִׂם אֶת הַלֻּחֹת בָּאָרוֹן אֲשֶׁר עָשִׂיתִי וַיִּהְיוּ שָׁם כַּאֲשֶׁר צִוַּנִי יְהֹוָה׃}
}
{
	\eDtr1{
		And I turned and went down from the mountain, and I set
		the tablets in the ark that I had made, and they have been
		there, as YHWH commanded me.
	}
}
{
}

\Verse{6}
{
	\hDtr1{וּבְנֵי יִשְׂרָאֵל נָסְעוּ מִבְּאֵרֹת בְּנֵי יַעֲקָן מוֹסֵרָה שָׁם מֵת אַהֲרֹן וַיִּקָּבֵר שָׁם וַיְכַהֵן אֶלְעָזָר בְּנוֹ תַּחְתָּיו׃}
}
{
	\eDtr1{
		“And the children of Israel had traveled from
		Beeroth-bene-Jaakan to Moserah. There Aaron died, and he
		was buried there; and Eleazar, his son, functioned as
		priest in his place.
	}
}
{
}

\Verse{7}
{
	\hDtr1{מִשָּׁם נָסְעוּ הַגֻּדְגֹּדָה וּמִן הַגֻּדְגֹּדָה יטְבָתָה אֶרֶץ נַחֲלֵי מָיִם׃}
}
{
	\eDtr1{
		From there they traveled to Gudgod, and from Gudgod to
		Jotbah, a land of wadis of water.
	}
}
{
}

\Verse{8}
{
	\hDtr1{בָּעֵת הַהִוא הִבְדִּיל יְהֹוָה אֶת שֵׁבֶט הַלֵּוִי לָשֵׂאת אֶת אֲרוֹן בְּרִית יְהֹוָה לַעֲמֹד לִפְנֵי יְהֹוָה לְשָׁרְתוֹ וּלְבָרֵךְ בִּשְׁמוֹ עַד הַיּוֹם הַזֶּה׃}
}
{
	\eDtr1{
		At that time YHWH distinguished the tribe of Levi to carry
		the ark of YHWH’s covenant, to stand in front of YHWH to
		serve Him, and to bless in His name to this day.
	}
}
{
}

\Verse{9}
{
	\hDtr1{עַל כֵּן לֹא הָיָה לְלֵוִי חֵלֶק וְנַחֲלָה עִם אֶחָיו יְהֹוָה הוּא נַחֲלָתוֹ כַּאֲשֶׁר דִּבֶּר יְהֹוָה אֱלֹהֶיךָ לוֹ׃}
}
{
	\eDtr1{
		Therefore Levi has not had a portion and a legacy with its
		brothers. YHWH: He is its legacy, as YHWH, your God, spoke
		to it.
	}
}
{
}

\Verse{10}
{
	\hDtr1{וְאָנֹכִי עָמַדְתִּי בָהָר כַּיָּמִים הָרִאשֹׁנִים אַרְבָּעִים יוֹם וְאַרְבָּעִים לָיְלָה וַיִּשְׁמַע יְהֹוָה אֵלַי גַּם בַּפַּעַם הַהִוא לֹא אָבָה יְהֹוָה הַשְׁחִיתֶךָ׃}
}
{
	\eDtr1{
		“And I: I stood in the mountain as in the first days:
		forty days and forty nights. And YHWH listened to me that
		time as well. YHWH was not willing to destroy you.
	}
}
{
%	\footnote{
%		And I. Moses now resumes his account of his second forty days on Horeb,
%		which had left off at 9:29. He had interrupted his story with mention of
%		the ark and other matters.
%	}
}

\Verse{11}
{
	\hDtr1{וַיֹּאמֶר יְהֹוָה אֵלַי קוּם לֵךְ לְמַסַּע לִפְנֵי הָעָם וְיָבֹאוּ וְיִירְשׁוּ אֶת הָאָרֶץ אֲשֶׁר נִשְׁבַּעְתִּי לַאֲבֹתָם לָתֵת לָהֶם׃}
}
{
	\eDtr1{
		And YHWH said to me, ‘Get up. Set out on the journey in
		front of the people, and they’ll come and take possession
		of the land that I swore to their fathers to give to
		them.’
	}
}
{
}

\Verse{12}
{
	\hDtr1{וְעַתָּה יִשְׂרָאֵל מָה יְהֹוָה אֱלֹהֶיךָ שֹׁאֵל מֵעִמָּךְ כִּי אִם לְיִרְאָה אֶת יְהֹוָה אֱלֹהֶיךָ לָלֶכֶת בְּכל דְּרָכָיו וּלְאַהֲבָה אֹתוֹ וְלַעֲבֹד אֶת יְהֹוָה אֱלֹהֶיךָ בְּכל לְבָבְךָ וּבְכל נַפְשֶׁךָ׃}
}
{
	\eDtr1{
		“And now, Israel, what is YHWH, your God, asking from you
		except to fear YHWH, your God, to go in all His ways, and
		to love Him and to serve YHWH, your God, with all your
		heart and all your soul,
	}
}
{
%	\footnote{
%		what is God asking from you. All the things that follow in this verse
%		have been mentioned already in Moses’ speech. They come now, therefore,
%		as a summary, a statement of the basic things that underlie everything
%		else: fear, love, deeds, service, complete commitment.
%	}
}

\Verse{13}
{
	\hDtr1{לִשְׁמֹר אֶת מִצְוֺת יְהֹוָה וְאֶת חֻקֹּתָיו אֲשֶׁר אָנֹכִי מְצַוְּךָ הַיּוֹם לְטוֹב לָךְ׃}
}
{
	\eDtr1{
		to observe YHWH’s commandments and His laws that I command
		you today to be good for you.
	}
}
{
}

\Verse{14}
{
	\hDtr1{הֵן לַיהֹוָה אֱלֹהֶיךָ הַשָּׁמַיִם וּשְׁמֵי הַשָּׁמָיִם הָאָרֶץ וְכל אֲשֶׁר בָּהּ׃}
}
{
	\eDtr1{
		Here, YHWH, your God, has the skies—and the skies of the
		skies!—the earth and everything that’s in it.
	}
}
{
%	\footnote{
%		the skies of the skies. Meaning, presumably, that if one were in the sky
%		and looked up, one would see a farther portion of the sky. So the phrase
%		means the farthest reaches of the sky.
%	}
}

\Verse{15}
{
	\hDtr1{רַק בַּאֲבֹתֶיךָ חָשַׁק יְהֹוָה לְאַהֲבָה אוֹתָם וַיִּבְחַר בְּזַרְעָם אַחֲרֵיהֶם בָּכֶם מִכּל הָעַמִּים כַּיּוֹם הַזֶּה׃}
}
{
	\eDtr1{
		Only, YHWH was attracted to your fathers, to love them,
		and He chose their seed after them: you, out of all the
		peoples, as it is this day.
	}
}
{
}

\Verse{16}
{
	\hDtr1{וּמַלְתֶּם אֵת ערְלַת לְבַבְכֶם וְערְפְּכֶם לֹא תַקְשׁוּ עוֹד׃}
}
{
	\eDtr1{
		So you shall circumcise the foreskin of your heart, and
		you shall not harden your necks anymore.
	}
}
{
%	\footnote{
%		the foreskin of your heart. This concept occurs in two books of the
%		Torah (Deut 10:16; 30:6; Lev 26:41) and in two of the prophets (Jer 4:4;
%		Ezek 44:7,9). The sign of the Abrahamic covenant, and probably the most
%		commonly observed of the commandments by Jewish families for centuries,
%		is circumcision. But the Torah commands the circumcision of one’s heart
%		as well. It establishes that outward fulfillment of practices without
%		also feeling it in one’s heart is insufficient. The concept of
%		circumcising one’s heart, moreover, unlike physical circumcision,
%		applies to both women and men.
%	}
}

\Verse{17}
{
	\hDtr1{כִּי יְהֹוָה אֱלֹהֵיכֶם הוּא אֱלֹהֵי הָאֱלֹהִים וַאֲדֹנֵי הָאֲדֹנִים הָאֵל הַגָּדֹל הַגִּבֹּר וְהַנּוֹרָא אֲשֶׁר לֹא יִשָּׂא פָנִים וְלֹא יִקַּח שֹׁחַד׃}
}
{
	\eDtr1{
		Because YHWH, your God: He is the God of gods and the Lord
		of lords, the great, the mighty, and the awesome God, who
		won’t be partial and won’t take a bribe,
	}
}
{
%	\footnote{
%		the God of gods and the Lord of lords. Admittedly, this sounds
%		unmonotheistic, seemingly acknowledging the existence of other gods.
%		But, earlier, Moses has said, “YHWH: He is God. There is no other
%		outside of Him” (4:35) and “YHWH is one” (6:4). And, later, YHWH says,
%		“There is no god with me” (32:39). How are we to reconcile the many
%		monotheistic passages in Deuteronomy with “God of gods and Lord of
%		lords”? As in the case of not having “other gods before my face,” this
%		is a linguistic matter. “God of gods” is an expression for conveying
%		that YHWH is something that all other gods in whom people believe are
%		not. It need not presume that such other gods exist. Footnote 10:17.
%		won’t take a bribe. What could it mean to bribe God?! A sacrifice, a
%		donation, a kindness, charity, prayer: one might do any of these to win
%		the deity’s favor, rather than doing them as fulfillment of the
%		commandments or as free acts. Moses informs the people that this will
%		not work.
%	}
}

\Verse{18}
{
	\hDtr1{עֹשֶׂה מִשְׁפַּט יָתוֹם וְאַלְמָנָה וְאֹהֵב גֵּר לָתֶת לוֹ לֶחֶם וְשִׂמְלָה׃}
}
{
	\eDtr1{
		doing judgment for an orphan and a widow and loving an
		alien, to give him bread and a garment.
	}
}
{
}

\Verse{19}
{
	\hDtr1{וַאֲהַבְתֶּם אֶת הַגֵּר כִּי גֵרִים הֱיִיתֶם בְּאֶרֶץ מִצְרָיִם׃}
}
{
	\eDtr1{
		So you shall love the alien, because you were aliens in
		the land of Egypt.
	}
}
{
}

\Verse{20}
{
	\hDtr1{אֶת יְהֹוָה אֱלֹהֶיךָ תִּירָא אֹתוֹ תַעֲבֹד וּבוֹ תִדְבָּק וּבִשְׁמוֹ תִּשָּׁבֵעַ׃}
}
{
	\eDtr1{
		You shall fear YHWH, your God, you shall serve Him, and
		you shall cling to Him, and you shall swear by His name.
	}
}
{
}

\Verse{21}
{
	\hDtr1{הוּא תְהִלָּתְךָ וְהוּא אֱלֹהֶיךָ אֲשֶׁר עָשָׂה אִתְּךָ אֶת הַגְּדֹלֹת וְאֶת הַנּוֹרָאֹת הָאֵלֶּה אֲשֶׁר רָאוּ עֵינֶיךָ׃}
}
{
	\eDtr1{
		He is your splendor, and He is your God, who did these
		great and awesome things for you that your eyes have seen.
		Your fathers went down to Egypt with seventy persons, and
		now YHWH, your God, has made you like the stars of the
		skies for multitude.
	}
}
{
}

\Verse{22}
{
	\hDtr1{בְּשִׁבְעִים נֶפֶשׁ יָרְדוּ אֲבֹתֶיךָ מִצְרָיְמָה וְעַתָּה שָׂמְךָ יְהֹוָה אֱלֹהֶיךָ כְּכוֹכְבֵי הַשָּׁמַיִם לָרֹב׃}
}
{
	\eDtr1{}
}
{
}
